% page 228 of the facsimile (image p-227 of the scan)
% block 160.0 x 242.0 mm ; left margin 8.0 mm
\pgc{-12.700mm}{0.000mm}{
\l{\cel{1}220}
\l{}
\l{\sou{his}terio. (\sou{patol.}) Morbo nervala quan karakterizas sufoki,}
\l{\cel{3}konvulsi, e c., judikita (nejuste) kom partikulara ye}
\l{\cel{3}la homini. - DEFIRS.}
\l{}
\l{\sou{histerezometro}. (\sou{elektrotekniko)}. Aparato per qua onu mezuras}
\l{\cel{3}la qualeso di la fero uzata elektrale, de la vidopunto}
\l{\cel{3}\textquotedbl{}histerezo\textquotedbl{}. - DEFIS.}
\l{}
\l{\sou{histerotomio}. (\sou{kirurgio}). Disseko di la utero. - DEF.}
\l{}
\l{\sou{histogenezo}. (\sou{biol.}) Formaceso e kresko di la tisui diversa}
\l{\cel{3}di la embriono, de la blastomeri komencala. - DEFIRS.}
\l{}
\l{\sou{histolizo.(biol.})\cel{2}Destrukto di la tisui vivanta, precipue en}
\l{\cel{3}la komenco di la metamorfosi.}
\l{}
\l{\sou{histologio}. La parto di anatomio qua studias la formaco, la}
\l{\cel{3}evoluciono e la kompozeso di la tisui vivanta. - DEFIRS.}
\l{}
\l{historio. I. Biografio di populo, di naciono. - II. Studio di}
\l{\cel{3}la enti diversa qui existas en naturo. - EFIRS.}
\l{}
\l{\sou{historiogra}fo. La persono qua komisesis oficale da princo,}
\l{\cel{3}pri kompozar la historio di la regno-duro di ica.- DEFIRS.}
\l{}
\l{\sou{historiografio}. La arto di la historiografo. - DEFIRS.}
\l{}
\l{histriono. I. (en la epoki antiqua). Aktoro qua pleis farsi}
\l{\cel{3}grosiera. - II. Bufono en la ferii. - III. Komedio-}
\l{\cel{3}pleistacho. - DEFIS.}
\l{}
\l{\sou{hivernar}. (\sou{netrans.}) Travivar la vintro-tempo shirmate. -EFIS.}
\l{}
\l{ho !\cel{3}I. Interjeciono per qua onu advokas . - II. Interjeciono}
\l{\cel{3}qua expresas sprico di joyo, di doloro, di timo. - III.}
\l{\cel{3}Interjeciono qua expresas surprizeso, advoko subita, sprico}
\l{\cel{3}da la psiko kande lu emocas forte. - IV. Interjeciono per}
\l{\cel{3}qua onu expresas indigneso od astoneso.}
\l{}
\l{\sou{+ho-}.\cel{2}Prefixo qua signifikas \textquotedbl{}(radiko) dum qua ni vivas nun\textquotedbl{}}
\l{\cel{3}(kom ex.: hodie, homonate, hosemane).}
\l{}
\l{\sou{+hobio.}\cel{2}Ideo quan onu preferas traktar, o pri qua onu preferas}
\l{\cel{3}informesar ; kozo pri qua onu su okupas prefere e kom dis-}
\l{\cel{3}traktivo, amuzivo. - E.}
\l{}
\l{\sou{hoboyo}. Muzik-instrumento vento-funcionanta, langetoza, sen}
\l{\cel{3}beko, kono-forma, qua finas per mikra funelo, e qua emisas}
\l{\cel{3}foni klara e tre dolca. - DEFIRS.}
\l{}
\l{\sou{hodografo.}\cel{2}(\sou{matem.})\cel{2}De punto fixa O, onu trasas singla-}
\l{\cel{3}instante vektoro qua equipolas la rapideso-vektoro. La ex-}
\l{\cel{3}tremajo di ta vektoro deskriptas kurvo qua esas la \textquotedbl{}hodo-}
\l{\cel{3}grafo\textquotedbl{} di ta movado.}
}
